% Appendix, from exam/README.md: what the two written papers are for, what they cover, the
% conventions they assume, and how they are marked.

\chapter{The Self-Assessment Papers}
\label{exam:ch:about}

This appendix and the four after it hold two complete four-hour papers for \emph{Embedded Linux and
Kernel Drivers}, with model answers. Appendices~\ref{exa:ch:paper} and \ref{exb:ch:paper} are the
papers, questions only; Appendices~\ref{ansa:ch:answers} and \ref{ansb:ch:answers} are their model
answers, with marks.

\section*{What these are for}
The book has \textbf{no exam and nothing is marked}. What you get instead is a lab per chapter,
checked on the target by \code{make test}, and eight to ten exercises per chapter ending in a
Cross-check that makes you compute a number by hand and then measure it.

These papers do not replace that. \textbf{They are here to test your own skills and knowledge,
nothing more.} They gate nothing, they are not a qualification, and no part of the book requires
them. Nothing in the companion repository depends on them, and \code{make test} does not know they
exist. They are useful where a written result is wanted anyway, for a certifying employer or a
formal course credit, and useful on their own for finding out what you can reconstruct with nothing
in front of you.

\textbf{Take one after the last chapter.} Both papers draw on all twelve chapters, so sitting one
partway through examines material you have not read yet, and the result says more about how far you
have read than about what you have understood. The intended point is after Chapter~\ref{c12:ch},
once the driver probes off a device tree, services interrupts, blocks readers correctly, registers
with two subsystems, and has had its latency characterised.

\section*{The two papers}
Both cover the whole book, Chapters~1 to 12, and share no question. Either can be used alone; use
both as a main sitting and a resit, or in alternate years. They do \textbf{not} weight the chapters
identically, and that is deliberate.

\subsection*{Paper A}
\begin{booktable}{@{}p{4.6em}Lp{5em}r@{}}
  \thead{Question} & \thead{Topic} & \thead{Chapters} & \thead{Marks} \\
  \midrule
  1 & The image, and the licence on it & 1 & 12 \\
  2 & Asking a running kernel a question & 2 & 8 \\
  3 & Configuring a kernel, and what it costs & 3 & 10 \\
  4 & A module, written out & 4 & 12 \\
  5 & The character device contract & 5 & 12 \\
  6 & An address, translated and mapped & 6, 10 & 12 \\
  7 & A race, traced & 7 & 10 \\
  8 & An interrupt, acknowledged & 8 & 10 \\
  9 & A reader put to sleep & 9 & 8 \\
  10 & Frameworks, and the trade & 11, 12 & 6 \\
  \midrule
    & & & \textbf{100} \\
\end{booktable}

\subsection*{Paper B}
\begin{booktable}{@{}p{4.6em}Lp{5em}r@{}}
  \thead{Question} & \thead{Topic} & \thead{Chapters} & \thead{Marks} \\
  \midrule
  1 & What leaves the building, and when it does not & 1 & 10 \\
  2 & A number that is not what it says & 2, 9 & 12 \\
  3 & An option that did not survive & 3 & 10 \\
  4 & A module that will not load & 4 & 12 \\
  5 & A read that returns zero & 5, 9 & 12 \\
  6 & A mapping that is wrong & 6 & 10 \\
  7 & A deadlock that has not happened & 7 & 12 \\
  8 & An interrupt that never stops & 8 & 10 \\
  9 & A driver that finds its own hardware & 10, 11 & 6 \\
  10 & Worse average, better maximum & 12 & 6 \\
  \midrule
    & & & \textbf{100} \\
\end{booktable}

\textbf{Paper A leans towards tracing a mechanism forwards}: an address translated through a device
tree, a module written from nothing, a \code{read()} traced through its four cases, an interrupt
acknowledged in the right order.

\textbf{Paper B leans towards the conditions under which the mechanism stops working}: an option
silently dropped by \code{olddefconfig}, a module refused at load, a read that returns end of file
when it meant ``not yet'', a mapping that reads plausible zeros, a deadlock reported by a tool
without the deadlock occurring.

\section*{Code on a written paper}
Both papers \textbf{do ask a candidate to write kernel C}. That is deliberate. This book teaches by
building a driver: the specification of every lab is prose in its chapter and a test suite that
tells you whether what you wrote satisfies it, and a paper that avoided code entirely would be
examining a different course.

The code questions come in two shapes, and neither is a typing exercise:
\begin{itemize}
  \item \textbf{Write one function.} An init function, a \code{read()}, an interrupt handler. The
    answers are ten to twenty lines each. What carries the marks is the algorithm: the error path,
    the return value, the accessor used, the context the code is running in. Not the syntax.
  \item \textbf{Find what is wrong.} Both papers hand over a listing that compiles and does
    something other than what its author intended. The defects are the ones the chapters single
    out.
\end{itemize}
Both papers state in their rubric that answers may be written in kernel C \textbf{or in unambiguous
pseudocode}, and that nothing is marked on syntax. A missing semicolon costs nothing. A
\code{copy_to_user} replaced by \code{memcpy} costs everything.

\section*{Conventions the papers assume}
Both papers state these in their own rubric, so a candidate never has to have read this appendix.
\begin{itemize}
  \item \textbf{The kernel is 6.12 as pinned by the course.} Where an interface has changed within
    recent memory the paper says which version it means, because Chapters~\ref{c3:ch}, \ref{c5:ch}
    and~\ref{c11:ch} all make the point that the internal API is not stable.
  \item \textbf{\code{HZ} is 250} unless a question says otherwise, so a jiffy is 4 ms.
  \item \textbf{The device is \code{qa-dev}}, whose register map is the one in the course, and any
    register a question depends on is quoted in the question.
  \item \textbf{Only the values nobody could be expected to carry are supplied}: the SPI base of 32,
    the page size, a register offset where one is needed. Address translations, jiffy arithmetic and
    latency budgets are derived, because deriving them is the examinable skill.
  \item \textbf{``State the consequence''} means naming the defect earns half the marks and saying
    what it does to the running system earns the other half. ``This is wrong'' scores nothing.
  \item \textbf{Where a question asks you to find defects, the number is stated.} Listing more is
    not penalised, but only the stated number is marked, so put your strongest answers first.
\end{itemize}

\section*{Marking}
Every model answer is written to be marked by somebody who has read the chapters and does not
otherwise write kernel code, so each carries the reasoning rather than the answer alone. Marks are
shown per part.

Three conventions worth agreeing before a paper is marked:
\begin{itemize}
  \item \textbf{Method carries the marks.} A correct approach with an arithmetic slip in it is worth
    more than a correct number with no working, and both papers are built so that later parts
    consume earlier answers. Follow through an error rather than penalising it twice. A candidate
    who translates a device tree address wrongly and then maps their own wrong address correctly has
    demonstrated the examinable skill.
  \item \textbf{The named traps are worth full marks on their own.} Several questions exist entirely
    to see whether a candidate avoids one specific mistake: \code{read} returning 0 when it meant
    ``not yet'', a \code{memcpy} where \code{copy_to_user} was needed, dividing \file{/proc/stat} by
    \code{HZ} rather than by \code{USER_HZ}, taking \code{reg} at face value without translating
    through \code{ranges}, concluding from a clean unlocked run that a lock is unnecessary. Where a
    model answer flags one of these, a candidate who walks into it loses those marks and no others.
  \item \textbf{The discussion parts are not decoration.} ``State the consequence'' and ``say what
    this is blind to'' are where the book's actual content is, and a paper marked only on the code
    and the arithmetic would pass a candidate who has understood none of it. Paper~B in particular
    is mostly these.
\end{itemize}

\textbf{The code questions are answered as checklists, not as listings.} Where a paper asks a
candidate to write a function, the model answer states what a correct answer must contain and what
each element is worth, rather than printing a reference implementation. That is deliberate twice
over: it is how these answers are actually marked, since two correct answers will not look alike;
and a worked listing here would be a published solution to the lab that asks for the same thing,
and the book publishes none for its labs or its Code exercises. Appendix~\ref{sol:ch:solutions}
answers only the exercises that are not code.

\textbf{Where a model answer says a candidate may reasonably disagree, they may.} Two questions ask
for a judgement rather than a fact: Paper~A's Question~10(c) and Paper~B's Question~8(d). Both
answers state a position and say what an answer arguing the other way must contain to earn the
marks.
